% !TEX program = pdflatex
\documentclass[11pt, a4paper]{article}

\usepackage{amsmath, amssymb}
\usepackage{booktabs}
\usepackage{geometry}
\usepackage{microtype}
\usepackage{parskip}
\usepackage{xcolor}
\usepackage{listings}
\usepackage{hyperref}
\usepackage{enumitem}
\usepackage{titlesec}
\geometry{margin=2.5cm}

% ── Python listing style ──────────────────────────────────────────────────────
\definecolor{codebg}{RGB}{248,248,248}
\definecolor{codeframe}{RGB}{200,200,200}
\definecolor{kw}{RGB}{0,0,180}
\definecolor{str}{RGB}{180,0,0}
\definecolor{cmt}{RGB}{80,120,80}
\definecolor{num}{RGB}{0,100,0}

\lstdefinestyle{py}{
  language        = Python,
  backgroundcolor = \color{codebg},
  frame           = single,
  rulecolor       = \color{codeframe},
  basicstyle      = \ttfamily\small,
  keywordstyle    = \color{kw}\bfseries,
  stringstyle     = \color{str},
  commentstyle    = \color{cmt}\itshape,
  numberstyle     = \tiny\color{gray},
  numbers         = left,
  numbersep       = 6pt,
  breaklines      = true,
  breakatwhitespace = true,
  showstringspaces= false,
  tabsize         = 4,
  captionpos      = b,
  morekeywords    = {self, True, False, None, lambda, as},
}
\lstset{style=py}

\hypersetup{colorlinks=true, linkcolor=blue!60!black, urlcolor=blue!60!black}

% ── Title ─────────────────────────────────────────────────────────────────────
\title{\textbf{v3 Codebase Explained}\\[4pt]
       \large Environment and Training Script — Line-by-Line}
\author{}
\date{}

\begin{document}
\maketitle
\tableofcontents
\newpage

% =============================================================================
\section{Overview}
% =============================================================================

The codebase consists of two Python files:

\begin{itemize}[noitemsep]
  \item \texttt{Environments/v3.py} --- the Gymnasium environment that simulates
        a sector of en-route airspace, computes rewards, and exposes observations
        to the RL agent.
  \item \texttt{Training/v3\_train.py} --- the PPO training script that wraps
        the environment, configures Stable-Baselines3, and manages checkpoints
        and logging.
\end{itemize}

The environment follows the standard Gymnasium interface: \texttt{reset()} starts
a new episode, \texttt{step(action)} advances by one RL step (5 simulated seconds),
and returns \texttt{(obs, reward, terminated, truncated, info)}.

% =============================================================================
\section{Environment: \texttt{Environments/v3.py}}
% =============================================================================

% -----------------------------------------------------------------------------
\subsection{Imports and dependencies}
% -----------------------------------------------------------------------------

\begin{lstlisting}[caption={Module-level imports.}]
import math, random
import numpy as np
import gymnasium as gym
from gymnasium import spaces
import bluesky as bs
from bluesky.simulation import ScreenIO
from bluesky.tools import geo
from bluesky.stack.stackbase import Stack as _BsStack
from polygenerator import random_convex_polygon
from shapely.geometry import Polygon as ShapelyPolygon
from shapely.affinity import scale as shapely_scale
\end{lstlisting}

\begin{itemize}[noitemsep]
  \item \textbf{bluesky} --- open-source ATC simulator; provides aircraft
        physics, position tracking (\texttt{bs.traf}), and a command stack
        (\texttt{bs.stack}).
  \item \textbf{geo} --- BlueSky's geodesic helpers: \texttt{kwikqdrdist}
        (quick bearing + distance), \texttt{kwikdist} (distance only),
        \texttt{qdrpos} (project a point along a bearing).
  \item \textbf{polygenerator} --- generates random convex polygons as vertex
        lists, used to build the randomised sector boundary.
  \item \textbf{shapely} --- computational geometry; used to scale the polygon
        to the required area and to interpolate points on the boundary.
\end{itemize}

% -----------------------------------------------------------------------------
\subsection{Configuration dictionary}
% -----------------------------------------------------------------------------

\begin{lstlisting}[caption={The \texttt{CONFIG} dictionary --- all tuneable parameters
in one place. Lambda values are re-sampled fresh each episode.}]
CONFIG = {
    # Aircraft & sector
    'ac_type':               'A320',
    'ac_speed':              450.0,        # kt TAS, fixed throughout episode
    'ac_mach':               0.78,
    'altitude':              350,          # FL350
    'center_ll':             (0.0, 0.0),   # flat-earth equatorial: cos(0)=1
    'n_aircraft':            lambda: random.randint(10, 15),
    'rho':                   lambda: random.uniform(1/15000, 1/5000),
    'sep_nm':                5.0,          # ICAO horizontal separation standard
    'buffer_nm':             10.0,         # extra clearance at spawn
    'dest_dist_factor':      2.0,          # destination = 2 x sector diameters away
    # Polygon
    'n_vertices':            lambda: random.randint(5, 7),
    'min_circularity':       0.65,         # reject very elongated polygons
    'max_placement_tries':   50,
    # Simulation
    'sim_dt':                0.5,          # BlueSky physics timestep (s)
    'action_freq':           10,           # RL step = 10 x 0.5 s = 5 s
    'lookahead_s':           900.0,        # 15 min conflict detection horizon
    't_warn':                600.0,        # 10 min warning horizon
    'crossings_per_episode': 4.0,
    # Focus selection
    'focus_clear_steps':     5,
    'focus_emergency_u':     0.8,
    # Reward weights
    'w_los':                 10.00,
    'w_conflict':            3.00,
    'w_drift':               0.60,
    'w_work':                0.50,
}
\end{lstlisting}

\paragraph{Density parameter $\rho$.}
The sector area is computed as $A = N / \rho$, where $N$ is the number of aircraft
and $\rho$ is the aircraft density in aircraft/km$^2$.
The range $[1/15\,000,\, 1/5\,000]$ aircraft/km$^2$ corresponds to
5\,000--15\,000\,km$^2$ per aircraft.
Defining $\rho$ rather than area directly means the sector automatically scales
with traffic count, keeping density constant across different $N$ draws.

% -----------------------------------------------------------------------------
\subsection{Global constants}
% -----------------------------------------------------------------------------

\begin{lstlisting}[caption={Module-level constants derived from \texttt{CONFIG}.}]
N_NBR   = CONFIG['n_neighbours']   # 4 intruder slots in the observation
OBS_DIM = 3 + N_NBR * 5           # 3 ownship + 4x5 intruder = 23 floats

D_WARN = CONFIG['t_warn'] * CONFIG['ac_speed'] / 3600.0  # 75 NM
V_NOM  = CONFIG['ac_speed'] / 3600.0                      # NM/s

TURN_DELTAS = {0: -60, 1: -45, 2: -30, 4: 30, 5: 45, 6: 60}
ACT_COST    = [1.0, 0.75, 0.5, 0.0, 0.5, 0.75, 1.0, None]
\end{lstlisting}

\texttt{D\_WARN} is the distance at which an aircraft moving at cruise speed
would reach the warning horizon: $d_\text{warn} = T_\text{warn} \times v = 600\,\text{s}
\times 450\,\text{kt} / 3600 = 75\,\text{NM}$.

\texttt{ACT\_COST} encodes the workload cost of each action.
The costs are proportional to $|\delta|/60^\circ$, so that the total cost depends
only on the accumulated heading excursion, not on how many instructions achieved it.
Action~3 (hold) is free; index~7 (fly direct) is \texttt{None} because its cost
is computed dynamically.

% -----------------------------------------------------------------------------
\subsection{Coordinate helpers}
% -----------------------------------------------------------------------------

\begin{lstlisting}[caption={Flat-earth projection between geographic and NM coordinates.}]
def latlon_to_nm(center_ll, lat, lon):
    """Convert (lat, lon) to (east_nm, north_nm) relative to center_ll."""
    ref_lat, ref_lon = center_ll
    east_nm  = (lon - ref_lon) * 60.0 * math.cos(math.radians(ref_lat))
    north_nm = (lat - ref_lat) * 60.0
    return np.array([east_nm, north_nm])

def nm_to_latlon(center_ll, east_nm, north_nm):
    """Convert (east_nm, north_nm) offsets back to (lat, lon)."""
    ref_lat, ref_lon = center_ll
    return (ref_lat + north_nm / 60.0,
            ref_lon + east_nm  / (60.0 * math.cos(math.radians(ref_lat))))

def wrap_to_180(angle_deg):
    """Wrap an angle in degrees to (-180, 180]."""
    return (angle_deg + 180) % 360 - 180
\end{lstlisting}

Because \texttt{center\_ll = (0.0, 0.0)}, $\cos(\phi_\text{ref}) = 1$ and the
longitude conversion simplifies to $\Delta\text{lon} \times 60$.
One degree of latitude or longitude thus equals exactly 60\,NM, giving a
pure flat-earth Cartesian projection.

All geometry in the codebase --- CPA computations, separation checks, urgency ---
works in NM east/north coordinates returned by \texttt{latlon\_to\_nm}.

% -----------------------------------------------------------------------------
\subsection{Sector polygon}
% -----------------------------------------------------------------------------

\begin{lstlisting}[caption={Generating a random convex sector polygon of the required area.}]
def _make_polygon(area_km2):
    target_nm2 = area_km2 * KM_TO_NM ** 2
    scaled = None
    for _ in range(1000):
        raw   = ShapelyPolygon(random_convex_polygon(CONFIG['n_vertices']()))
        scale = math.sqrt(target_nm2 / raw.area)
        scaled = shapely_scale(raw, xfact=scale, yfact=scale, origin='centroid')
        if 4 * math.pi * scaled.area / scaled.length ** 2 >= CONFIG['min_circularity']:
            break
    cx, cy = scaled.centroid.x, scaled.centroid.y
    return ShapelyPolygon([(x - cx, y - cy) for x, y in scaled.exterior.coords])
\end{lstlisting}

The function:
\begin{enumerate}[noitemsep]
  \item Samples a random unit convex polygon using \texttt{polygenerator}.
  \item Scales it uniformly so its area equals the target (in NM$^2$).
  \item Rejects the polygon if its isoperimetric ratio $4\pi A / L^2 < 0.65$,
        which would indicate a very elongated or spiky shape that would produce
        unrealistic flight geometries.
  \item Centres the accepted polygon at the origin by subtracting the centroid.
\end{enumerate}

% -----------------------------------------------------------------------------
\subsection{Aircraft placement: \texttt{\_place\_one}}
% -----------------------------------------------------------------------------

\begin{lstlisting}[caption={Placing one aircraft on the sector boundary.}]
def _place_one(polygon, sector, n_sectors):
    minx, miny, maxx, maxy = polygon.bounds
    dest_dist = math.sqrt((maxx-minx)**2 + (maxy-miny)**2) * CONFIG['dest_dist_factor']

    t_spawn  = (sector + CONFIG['spawn_jitter']()) / n_sectors
    t_ref    = (t_spawn + 0.5 + CONFIG['ref_jitter']()) % 1.0
    spawn_pt = polygon.exterior.interpolate(t_spawn, normalized=True)
    ref_pt   = polygon.exterior.interpolate(t_ref,   normalized=True)

    spawn_ll    = nm_to_latlon(CONFIG['center_ll'], spawn_pt.x, spawn_pt.y)
    ref_ll      = nm_to_latlon(CONFIG['center_ll'], ref_pt.x,   ref_pt.y)
    spawn_hdg, _ = geo.kwikqdrdist(*[float(v) for v in (*spawn_ll, *ref_ll)])

    dest_lat, dest_lon = geo.qdrpos(
        float(spawn_ll[0]), float(spawn_ll[1]), spawn_hdg, dest_dist)

    return {'sp_ll': spawn_ll, 'dest_ll': (dest_lat, dest_lon), 'heading': spawn_hdg}
\end{lstlisting}

The sector perimeter is parameterised as a normalised arc length $t \in [0,1)$.
Aircraft $k$ of $N$ is assigned $t_\text{spawn} = (k + \epsilon_1) / N$ where
$\epsilon_1 \sim \mathcal{U}(0.1, 0.9)$, distributing spawns evenly around the
boundary while avoiding exact symmetry.
The reference point that sets the initial heading is placed half-way around the
perimeter ($t_\text{ref} = t_\text{spawn} + 0.5$) with a small jitter, so the
aircraft initially points roughly across the sector.
The destination is placed two sector diameters beyond the reference bearing,
ensuring aircraft do not U-turn back into the sector.

% -----------------------------------------------------------------------------
\subsection{Pair urgency: \texttt{\_pair\_urgency}}
% -----------------------------------------------------------------------------

\begin{lstlisting}[caption={Computing the urgency $u_{ij}$ between two aircraft.}]
def _pair_urgency(i, j):
    pos_i  = latlon_to_nm(CONFIG['center_ll'], bs.traf.lat[i], bs.traf.lon[i])
    pos_j  = latlon_to_nm(CONFIG['center_ll'], bs.traf.lat[j], bs.traf.lon[j])
    d_east  = pos_j[0] - pos_i[0]
    d_north = pos_j[1] - pos_i[1]
    dist_sq = d_east**2 + d_north**2
    sep     = CONFIG['sep_nm']

    # Branch 1: active loss of separation
    if dist_sq < sep**2:
        return 1.0 + 9.0 * (1.0 - math.sqrt(dist_sq) / sep)

    # Relative velocity
    ve_i = spd_i * math.sin(math.radians(bs.traf.hdg[i]))
    vn_i = spd_i * math.cos(math.radians(bs.traf.hdg[i]))
    ve_j = spd_j * math.sin(math.radians(bs.traf.hdg[j]))
    vn_j = spd_j * math.cos(math.radians(bs.traf.hdg[j]))
    dv_east  = ve_j - ve_i
    dv_north = vn_j - vn_i

    rel_spd_sq = dv_east**2 + dv_north**2
    if rel_spd_sq < 1e-12:
        return 0.0   # same velocity vector, never converge

    # TCPA: time to closest point of approach
    range_rate = d_east * dv_east + d_north * dv_north   # r.v; negative = converging
    tcpa = -range_rate / rel_spd_sq
    if tcpa < 0 or tcpa > CONFIG['lookahead_s']:
        return 0.0   # diverging or beyond horizon

    # DCPA: miss distance at CPA
    dcpa_sq = max(0.0, dist_sq - range_rate**2 / rel_spd_sq)
    if dcpa_sq >= sep**2:
        return 0.0   # passes with safe margin

    # Branch 2: predicted conflict
    return min(1.0, max(0.0, (CONFIG['t_warn'] - tcpa) / CONFIG['t_warn']))
\end{lstlisting}

The urgency function implements the piecewise formula

\begin{equation}
u_{ij} =
\begin{cases}
1 + 9\!\left(1 - \dfrac{d_{ij}}{d_\text{sep}}\right)
  & d_{ij} < d_\text{sep} \quad\text{(active LoS)} \\[8pt]
\dfrac{T_\text{warn} - t_\text{CPA}}{T_\text{warn}}
  & d_\text{CPA} < d_\text{sep},\; 0 \le t_\text{CPA} \le T_\text{look} \\[6pt]
0 & \text{otherwise.}
\end{cases}
\end{equation}

\paragraph{CPA geometry.}
The time to closest approach is $t_\text{CPA} = -(\mathbf{r} \cdot \mathbf{v}) /
|\mathbf{v}|^2$, where $\mathbf{r} = \mathbf{p}_j - \mathbf{p}_i$ is the relative
position and $\mathbf{v} = \mathbf{v}_j - \mathbf{v}_i$ is the relative velocity.
The miss distance at CPA is
$d_\text{CPA}^2 = |\mathbf{r}|^2 - (\mathbf{r} \cdot \mathbf{v})^2 / |\mathbf{v}|^2$.
Note that $\mathbf{r} \cdot \mathbf{v} < 0$ (negative \texttt{range\_rate}) means
the aircraft are converging, giving $t_\text{CPA} > 0$.

% =============================================================================
\section{The \texttt{AirspaceEnv} class}
% =============================================================================

% -----------------------------------------------------------------------------
\subsection{State variables (\texttt{\_\_init\_\_})}
% -----------------------------------------------------------------------------

\begin{lstlisting}[caption={Instance variables initialised in \texttt{\_\_init\_\_}.}]
self.n_aircraft           = 0       # aircraft count for this episode
self._slots               = []      # list of length n_ac; each entry is a callsign or None
self._active_callsigns    = set()   # callsigns currently in the simulator
self._destination_ll      = {}      # cs -> (dest_lat, dest_lon)
self._commanded_heading   = {}      # cs -> current psi_cmd (deg)
self._steps_since_urgency = {}      # cs -> steps since any conflict was active
self._next_callsign_id    = 0       # monotonically increasing ID: AC00, AC01, ...
self._step_count          = 0
self._max_steps           = 0
self._focus_cs            = None    # callsign that receives the next instruction
self._pending_spawns      = {}      # slot -> countdown steps until respawn
self._urgency_matrix      = np.zeros((0, 0))   # N x N urgency matrix
self._urgency_cs_list     = []      # ordered list matching matrix rows/columns
self._los_this_step       = False   # set True if any pair violates sep this step
\end{lstlisting}

The \textit{slot} concept is important: there are always exactly \texttt{n\_ac}
slots, each either occupied by a callsign or \texttt{None}.
When an aircraft exits the sector, its slot is freed and a replacement is queued
in \texttt{\_pending\_spawns}.
This keeps the total aircraft count constant throughout the episode.

% -----------------------------------------------------------------------------
\subsection{Episode reset (\texttt{reset})}
% -----------------------------------------------------------------------------

\begin{lstlisting}[caption={Key steps in \texttt{reset()}.}]
def reset(self, seed=None, options=None):
    # 1. Seed Python random and NumPy for reproducibility
    effective_seed = seed if seed is not None else CONFIG['seed']
    ...

    # 2. Clear BlueSky state completely
    _BsStack.cmdstack.clear()
    bs.traf.reset()
    bs.stack.stack(f"DT {CONFIG['sim_dt']};FF")   # set timestep; fast-forward

    # 3. Sample episode parameters
    n_ac     = CONFIG['n_aircraft']()   # e.g. 12
    rho      = CONFIG['rho']()          # e.g. 1.2e-4 aircraft/km^2
    area_km2 = float(n_ac / rho)        # e.g. 100,000 km^2

    # 4. Build random sector polygon
    poly = _make_polygon(area_km2)
    ...

    # 5. Compute episode length from sector diameter and crossing count
    sector_diam_nm  = math.sqrt((maxx-minx)**2 + (maxy-miny)**2)
    step_duration_s = CONFIG['action_freq'] * CONFIG['sim_dt']   # 5 s
    self._max_steps = max(50, round(
        CONFIG['crossings_per_episode'] * sector_diam_nm / CONFIG['ac_speed']
        * 3600 / step_duration_s
    ))

    # 6. Register sector in BlueSky area filter; disable ASAS
    bs.tools.areafilter.defineArea('SECTOR', 'POLY', self._flat_latlon())
    bs.stack.stack('ASAS OFF')

    # 7. Spawn n_ac aircraft with minimum separation
    min_spawn_sep = CONFIG['sep_nm'] + CONFIG['buffer_nm']   # 15 NM
    for slot in range(n_ac):
        for _ in range(CONFIG['max_placement_tries']):
            ac = _place_one(self._polygon_shape, slot, n_ac)
            if clear_of_all:   # >= 15 NM from every already-spawned aircraft
                self._spawn_aircraft(slot, ac)
                break

    # 8. Select initial focus aircraft and return first observation
    self._focus_cs = self._select_focus_aircraft()
    return self._get_observation(), {}
\end{lstlisting}

The episode length formula ensures an aircraft crossing the sector at cruise speed
takes $\texttt{crossings\_per\_episode} = 4$ crossings worth of steps.
For a typical sector diameter of 300\,NM at 450\,kt this gives
$4 \times 300 / 450 \times 3600 / 5 \approx 1920$ steps $= 9600$\,s of simulation.

% -----------------------------------------------------------------------------
\subsection{RL step (\texttt{step})}
% -----------------------------------------------------------------------------

\begin{lstlisting}[caption={The \texttt{step()} method --- executed once per RL decision.}]
def step(self, action):
    # 1. Process any aircraft waiting to re-enter the sector
    self._process_pending_spawns()

    # 2. Record focus aircraft and its current commanded heading
    acting_cs   = self._focus_cs
    pre_cmd_hdg = self._commanded_heading.get(acting_cs) if acting_cs else None

    # 3. Apply the chosen action to the focus aircraft
    if acting_cs:
        self._apply_action(acting_cs, int(action))

    # 4. Advance simulation by action_freq = 10 sub-steps (5 simulated seconds)
    self._los_this_step = False
    for _ in range(CONFIG['action_freq']):
        bs.sim.step()
    self._check_los_now()   # geometric LoS check on final positions
    self._step_count += 1

    # 5. Remove aircraft that left the sector, queue replacements
    self._process_exits()

    # 6. Recompute urgency matrix, select new focus aircraft
    self._focus_cs = self._select_focus_aircraft()

    # 7. Compute reward (for the acting_cs at its post-step position)
    reward = self._compute_reward(acting_cs, int(action), pre_cmd_hdg)

    # 8. Update episode statistics
    if self._los_this_step:
        self._ep_stats['los'] += 1

    # 9. Build and return output tuple
    truncated = self._step_count >= self._max_steps
    return self._get_observation(), reward, False, truncated, info
\end{lstlisting}

The episode always ends by truncation (\texttt{truncated=True}); the
\texttt{terminated} flag is always \texttt{False} because there is no
terminal state --- the environment continues indefinitely within an episode.

The order of operations inside \texttt{step} is significant:
\texttt{\_check\_los\_now} runs \emph{before} \texttt{\_process\_exits} so that
LoS between a departing aircraft and others is still captured, even if the
departing aircraft is removed in the next line.

% -----------------------------------------------------------------------------
\subsection{Focus selection}
% -----------------------------------------------------------------------------

\subsubsection{\texttt{\_select\_focus\_aircraft}}

\begin{lstlisting}[caption={Selecting which aircraft receives the next instruction.}]
def _select_focus_aircraft(self):
    active      = [cs for cs in self._active_callsigns if bs.traf.id2idx(cs) >= 0]
    n_active    = len(active)
    sim_indices = [bs.traf.id2idx(cs) for cs in active]

    # Build N x N urgency matrix
    urgency_mat = np.zeros((n_active, n_active))
    for ii in range(n_active):
        for jj in range(ii + 1, n_active):
            u = _pair_urgency(sim_indices[ii], sim_indices[jj])
            urgency_mat[ii, jj] = urgency_mat[jj, ii] = u

    pair_max   = urgency_mat.max(axis=1)   # worst single conflict per aircraft
    total_load = urgency_mat.sum(axis=1)   # total burden across all pairs

    # Track how long each aircraft has been conflict-free
    for i, cs in enumerate(active):
        if pair_max[i] > 0:
            self._steps_since_urgency[cs] = 0
        else:
            self._steps_since_urgency[cs] += 1

    # Select: highest total burden, or most-drifted if no conflicts
    best_cs = (active[int(np.argmax(total_load))]
               if total_load.max() > 0 else self._drift_fallback(active))

    # Hysteresis: keep current focus unless resolved or emergency
    if self._focus_cs in active and best_cs != self._focus_cs:
        focus_pair_max = pair_max[active.index(self._focus_cs)]
        focus_resolved = self._steps_since_urgency[self._focus_cs] >= clear_steps
        emergency      = pair_max.max() >= CONFIG['focus_emergency_u']
        if (focus_pair_max > 0 or not focus_resolved) and not emergency:
            return self._focus_cs   # stay locked on current focus

    return best_cs
\end{lstlisting}

The selection criterion $f = \arg\max_i \sum_{j \ne i} u_{ij}$ prefers the
aircraft most burdened across \emph{multiple} moderate conflicts over one
involved in a single severe conflict --- important at high traffic densities.

The hysteresis lock prevents the focus from switching away mid-resolution:
as long as the current focus aircraft still has \emph{any} active urgency,
or fewer than \texttt{focus\_clear\_steps = 5} steps have elapsed since its
urgency cleared, the lock holds.
The emergency threshold $u_\text{emerg} = 0.8$ (approximately 120\,s to
predicted conflict) overrides the lock.

\subsubsection{\texttt{\_drift\_fallback}}

\begin{lstlisting}[caption={Fallback selection when no conflicts are active.}]
def _drift_fallback(self, active):
    for cs in sorted(active):
        dest_bearing, _ = geo.kwikqdrdist(...)
        hdg_err = wrap_to_180(dest_bearing - self._commanded_heading[cs])
        drift   = (1 - math.cos(math.radians(hdg_err))) / 2
        if drift > max_drift:
            max_drift, best_cs = drift, cs
    # hysteresis: only switch if best exceeds focus by more than margin
    if max_drift <= focus_drift + CONFIG['drift_switch_margin']:
        return self._focus_cs
    return best_cs
\end{lstlisting}

When no pair has urgency $>0$, focus falls to the aircraft with the largest
commanded-heading deviation $\delta = (1 - \cos\Delta\psi)/2 \in [0,1]$.
A margin of 0.05 prevents rapid switching between aircraft with similar drift.

% -----------------------------------------------------------------------------
\subsection{Reward computation}
% -----------------------------------------------------------------------------

\subsubsection{\texttt{\_compute\_reward}}

\begin{lstlisting}[caption={Reward function --- four additive negative terms.}]
def _compute_reward(self, acting_cs, action_idx, pre_cmd_hdg=None):

    # LoS penalty: heavy binary, sector-wide
    r_los = -CONFIG['w_los'] if self._los_this_step else 0.0

    # Conflict penalty: worst predicted collision course
    r_conflict = -CONFIG['w_conflict'] * self._conflict_score(acting_cs)

    # Drift penalty: per-step cost for off-route commanded heading
    if acting_cs and acting_cs in self._destination_ll:
        idx = bs.traf.id2idx(acting_cs)
        if idx >= 0:
            dest_bearing, _ = geo.kwikqdrdist(...)
            hdg_err = wrap_to_180(dest_bearing - cmd_hdg)
            r_drift = -CONFIG['w_drift'] * (1 - math.cos(math.radians(hdg_err))) / 2

    # Workload penalty: one-time cost per instruction
    if action_idx == 7 and pre_cmd_hdg is not None:
        deviation = abs(wrap_to_180(dest_bearing - pre_cmd_hdg))
        r_work = -CONFIG['w_work'] * min(deviation / 60.0, 1.0)
    else:
        r_work = -CONFIG['w_work'] * ACT_COST[action_idx]

    return float(r_los + r_conflict + r_drift + r_work)
\end{lstlisting}

The four components and their design rationale:

\begin{description}[noitemsep, leftmargin=0pt]
  \item[$r_\text{LoS}$] Binary, sector-wide: any pair violating 5\,NM costs
    $-w_\text{LoS} = -10$ regardless of the focus aircraft.
  \item[$r_\text{conflict}$] Graduated: scales with both imminence and severity
    of the \emph{worst} predicted conflict (see Section~\ref{sec:conflict_score}).
  \item[$r_\text{drift}$] Charged every step the focus aircraft is held off-route,
    creating a sustained pressure to execute fly-direct once the conflict is resolved.
    Uses $(1-\cos\Delta\psi)/2$ so it is exactly zero on-route and 1 at 180$^\circ$.
  \item[$r_\text{work}$] For turns (0--2, 4--6): cost $= |\delta|/60^\circ$ (fixed,
    proportional to heading change).
    For fly-direct (7): cost $= \min(|\psi_\text{dest} - \psi_\text{cmd}^-| / 60^\circ, 1)$,
    charging more for aborting a large turn and less for completing one.
\end{description}

\subsubsection{\texttt{\_conflict\_score}}
\label{sec:conflict_score}

\begin{lstlisting}[caption={Computing the worst conflict score across all intruders.}]
def _conflict_score(self, cs):
    worst_score = 0.0
    for other_cs in self._active_callsigns:
        ...
        if dist_nm < sep:
            worst_score = 1.0   # active LoS -> maximum score
            continue

        range_rate = d_east * dv_east + d_north * dv_north
        tcpa       = -range_rate / rel_spd_sq
        if tcpa < 0 or tcpa > CONFIG['lookahead_s']:
            continue

        cpa_east  = d_east  + tcpa * dv_east
        cpa_north = d_north + tcpa * dv_north
        if cpa_east**2 + cpa_north**2 >= sep**2:
            continue   # safe miss

        dcpa  = math.sqrt(cpa_east**2 + cpa_north**2)
        score = max(0.0, 1 - tcpa / t_warn) * max(0.0, 1 - dcpa / sep)
        worst_score = max(worst_score, score)
    return worst_score
\end{lstlisting}

The score $s_k = (1 - t_\text{CPA}/T_\text{warn}) \times (1 - d_\text{CPA}/d_\text{sep})$
is the product of two factors, both in $[0,1]$:
the first rises as the conflict becomes more imminent;
the second rises as the miss distance shrinks.
Only pairs with $d_\text{CPA} < d_\text{sep}$ (a true collision course) contribute.

% -----------------------------------------------------------------------------
\subsection{Loss-of-separation check: \texttt{\_check\_los\_now}}
% -----------------------------------------------------------------------------

\begin{lstlisting}[caption={Geometric LoS check run after each 5-second simulation window.}]
def _check_los_now(self):
    active  = [cs for cs in self._active_callsigns if bs.traf.id2idx(cs) >= 0]
    sep_sq  = CONFIG['sep_nm'] ** 2

    for ii in range(len(active)):
        pos_i = latlon_to_nm(CONFIG['center_ll'], ...)
        for jj in range(ii + 1, len(active)):
            pos_j = latlon_to_nm(CONFIG['center_ll'], ...)
            d_east  = pos_j[0] - pos_i[0]
            d_north = pos_j[1] - pos_i[1]
            if d_east**2 + d_north**2 < sep_sq:
                self._los_this_step = True
                return   # early exit on first violation
\end{lstlisting}

This is a brute-force $O(N^2)$ pairwise check.
It runs on the positions after all 10 simulation sub-steps, setting the flag that
drives the LoS penalty in \texttt{\_compute\_reward}.
It returns immediately on the first violation since only the binary flag matters
for the reward.

% -----------------------------------------------------------------------------
\subsection{Action application: \texttt{\_apply\_action}}
% -----------------------------------------------------------------------------

\begin{lstlisting}[caption={Applying a discrete action to the focus aircraft.}]
def _apply_action(self, cs, action_idx):
    idx = bs.traf.id2idx(cs)

    if action_idx in TURN_DELTAS:
        # Consecutive: add delta to current commanded heading (stackable)
        current_cmd = self._commanded_heading.get(cs, bs.traf.hdg[idx])
        self._commanded_heading[cs] = (current_cmd + TURN_DELTAS[action_idx]) % 360

    elif action_idx == 7:
        # Fly direct: reset to current bearing to destination
        dest_bearing, _ = geo.kwikqdrdist(bs.traf.lat[idx], bs.traf.lon[idx], ...)
        self._commanded_heading[cs] = float(dest_bearing) % 360

    # Action 3 = hold: psi_cmd unchanged; just re-issue to BlueSky
    bs.stack.stack(f'HDG {cs} {self._commanded_heading[cs]:.1f}')
\end{lstlisting}

The key design: turn instructions use \emph{consecutive} semantics ---
$\psi_\text{cmd} \leftarrow \psi_\text{cmd} + \delta$.
This differs from the earlier one-shot design ($\psi_\text{cmd} \leftarrow
\psi_\text{dest} + \delta$) by allowing turns to stack.
Two consecutive $-30^\circ$ turns produce $-60^\circ$ total deviation.

The \texttt{HDG} command sends the heading instruction to the BlueSky simulator,
which then smoothly turns the aircraft over the next simulation sub-steps.

% -----------------------------------------------------------------------------
\subsection{Observation construction: \texttt{\_get\_observation}}
% -----------------------------------------------------------------------------

\begin{lstlisting}[caption={Building the 23-element observation vector.}]
def _get_observation(self):
    # --- Ownship features (3) ---
    hdg_err       = wrap_to_180(dest_bearing - cmd_hdg)
    turn_progress = wrap_to_180(cmd_hdg - own_hdg) / 180.0   # [-1, 1]

    obs = [
        math.sin(math.radians(hdg_err)),   # sin/cos avoids +-180 discontinuity
        math.cos(math.radians(hdg_err)),
        turn_progress,
    ]

    # --- 4 intruder slots x 5 features each (20) ---
    for other_cs in self._active_callsigns:
        ...
        # Rotate displacement to ego heading frame
        ego_lat = d_east * cos_hdg - d_north * sin_hdg   # lateral (right)
        ego_fwd = d_east * sin_hdg + d_north * cos_hdg   # forward

        r_norm     = min(1.0, dist_nm / D_WARN)
        theta_norm = math.atan2(ego_lat, ego_fwd) / math.pi

        # CPA polar features
        if dist_nm < sep:
            tcpa_n, cpa_r_n, cpa_theta_n = 0.0, dist_nm/sep, atan2(ego_lat,ego_fwd)/pi
        else:
            # compute TCPA, CPA position, rotate to ego frame ...
            tcpa_n      = tcpa / t_warn
            cpa_r_n     = dcpa / sep
            cpa_theta_n = atan2(cpa_ego_lat, cpa_ego_fwd) / pi

        intruders.append((urgency, dist_nm, r_norm, theta_norm,
                          cpa_r_n, cpa_theta_n, tcpa_n, other_cs))

    # --- Fill slots: urgent first, then nearest ---
    urgent  = sorted([r for r in intruders if r[0] > 0], key=lambda r: -r[0])
    nearest = sorted(intruders, key=lambda r: r[1])
    # merge, keeping only N_NBR = 4 unique entries
    ...
    for slot_k in range(N_NBR):
        if slot_k < len(selected):
            obs += [r_n, theta_n, cpa_r, cpa_th, t_cpa]
        else:
            obs += [1.0, 0.0, 0.0, 0.0, 1.0]   # empty slot sentinel
\end{lstlisting}

\paragraph{Ego heading frame.}
All intruder positions and CPA offsets are rotated from geographic (east, north)
to the focus aircraft's heading frame (forward, lateral):
\begin{align}
  \text{ego\_lat} &= \Delta e \cos\psi - \Delta n \sin\psi & (\text{lateral / right}) \\
  \text{ego\_fwd} &= \Delta e \sin\psi + \Delta n \cos\psi & (\text{forward})
\end{align}
where $\psi$ is the focus aircraft's actual heading, $\Delta e$ and $\Delta n$ are
east and north displacements.

\paragraph{Polar encoding.}
Each intruder position is expressed as
$(r, \theta)$ where $r = d / d_\text{warn} \in [0,1]$ and
$\theta = \operatorname{atan2}(\text{ego\_lat}, \text{ego\_fwd}) / \pi \in [-1,1]$.
A value of $\theta = 0$ means directly ahead; $\theta = \pm 1$ means directly behind.
This encoding is rotation-equivariant within the ego frame and preserves the
sign of the lateral offset (negative = left, positive = right).

\paragraph{Slot filling.}
The four slots are filled by first inserting aircraft with active urgency
($u_{fj} > 0$) in descending urgency order, then filling remaining slots with the
nearest aircraft by current distance.
Slots with no real aircraft are padded with the sentinel $(1, 0, 0, 0, 1)$:
maximum range, no bearing, no CPA conflict, tcpa at horizon.

\paragraph{Observation table.}

\begin{table}[h]
\centering
\caption{Full 23-element observation vector.}
\label{tab:obs}
\smallskip
\begin{tabular}{clll}
\toprule
\textbf{Index} & \textbf{Name} & \textbf{Definition} & \textbf{Range} \\
\midrule
0 & $\sin\Delta\psi$ & $\sin(\psi_\text{dest} - \psi_\text{cmd})$ & $[-1,1]$ \\
1 & $\cos\Delta\psi$ & $\cos(\psi_\text{dest} - \psi_\text{cmd})$ & $[-1,1]$ \\
2 & $p_\text{turn}$  & $(\psi_\text{cmd} - \psi_\text{act}) / 180^\circ$ & $[-1,1]$ \\
\midrule
$3{+}5k$ & $r$         & $d_{fj} / d_\text{warn}$ & $[0,1]$ \\
$4{+}5k$ & $\theta$    & $\operatorname{atan2}(\text{ego\_lat}, \text{ego\_fwd})/\pi$ & $[-1,1]$ \\
$5{+}5k$ & $r^\text{CPA}$ & $d^\text{CPA}_{fj} / d_\text{sep}$ & $[0,1]$ \\
$6{+}5k$ & $\theta^\text{CPA}$ & CPA bearing in ego frame / $\pi$ & $[-1,1]$ \\
$7{+}5k$ & $\hat{t}_\text{CPA}$ & $t_\text{CPA} / T_\text{warn}$ & $[0,1]$ \\
\bottomrule
\multicolumn{4}{l}{\small $k \in \{0,1,2,3\}$ indexes the four intruder slots.}
\end{tabular}
\end{table}

% -----------------------------------------------------------------------------
\subsection{Traffic management: exits and spawning}
% -----------------------------------------------------------------------------

\subsubsection{\texttt{\_process\_exits} and \texttt{\_find\_exited}}

\begin{lstlisting}[caption={Removing aircraft that have left the sector.}]
def _find_exited(self):
    exited = []
    for cs in list(self._active_callsigns):
        idx = bs.traf.id2idx(cs)
        if idx < 0:              # aircraft deleted from BlueSky
            exited.append(cs); continue
        inside = bs.tools.areafilter.checkInside('SECTOR', ...)
        if not inside[0]:        # aircraft outside sector polygon
            exited.append(cs)
    return exited

def _process_exits(self):
    for cs in self._find_exited():
        slot = self._slots.index(cs)
        bs.traf.delete(bs.traf.id2idx(cs))
        self._active_callsigns.discard(cs)
        self._slots[slot] = None
        # clean up all per-aircraft dictionaries
        self._destination_ll.pop(cs, None)
        self._commanded_heading.pop(cs, None)
        self._steps_since_urgency.pop(cs, None)
        # queue slot for respawn
        if slot not in self._pending_spawns:
            self._pending_spawns[slot] = random.randint(*self._spawn_delay_range)
\end{lstlisting}

\subsubsection{\texttt{\_process\_pending\_spawns}}

\begin{lstlisting}[caption={Countdown-based respawn logic.}]
def _process_pending_spawns(self):
    ready    = sorted(slot for slot, t in self._pending_spawns.items() if t <= 1)
    requeued = set()

    for slot in ready:
        del self._pending_spawns[slot]
        ac = self._generate_replacement(slot)
        if ac is not None:
            self._spawn_aircraft(slot, ac)
        else:
            self._pending_spawns[slot] = 5   # placement failed: retry in 5 steps
            requeued.add(slot)

    for slot in list(self._pending_spawns):
        if slot not in requeued:
            self._pending_spawns[slot] -= 1   # decrement countdown
\end{lstlisting}

Each pending slot carries an integer countdown.
When the countdown reaches 1 the slot fires: a placement is attempted
(\texttt{\_generate\_replacement} tries up to 50 random positions, requiring
15\,NM clearance from all active aircraft).
If placement fails, the slot is re-queued with a 5-step retry delay.
All other pending slots are decremented.

% =============================================================================
\section{Training script: \texttt{Training/v3\_train.py}}
% =============================================================================

% -----------------------------------------------------------------------------
\subsection{Settings and PPO configuration}
% -----------------------------------------------------------------------------

\begin{lstlisting}[caption={Top-level training constants and PPO hyperparameters.}]
N_ENVS           = 4
TOTAL_TIMESTEPS  = 100_000_000   # 100 M steps
CHECKPOINT_EVERY = 100_000
N_RUNS           = 3             # for --multi mode

PPO_KWARGS = dict(
    learning_rate = 3e-4,
    n_steps       = 1024,     # rollout steps per environment before update
    batch_size    = 64,
    n_epochs      = 10,
    gamma         = 0.99,
    gae_lambda    = 0.95,
    clip_range    = 0.2,
    ent_coef      = 0.02,     # entropy bonus to encourage exploration
    vf_coef       = 0.5,
    policy_kwargs = dict(net_arch=[128, 128]),   # two hidden layers of 128 units
)
\end{lstlisting}

\begin{description}[noitemsep]
  \item[\texttt{N\_ENVS = 4}] Four environments run in parallel via
    \texttt{SubprocVecEnv}, giving a total rollout buffer of
    $4 \times 1024 = 4096$ transitions per update.
  \item[\texttt{n\_steps = 1024}] With 5\,s per step and 4 envs, each update
    uses $4 \times 1024 \times 5 = 20\,480$ seconds of simulated time.
  \item[\texttt{ent\_coef = 0.02}] The entropy coefficient encourages the policy
    to remain exploratory, reducing the risk of premature convergence to
    a degenerate hold-only strategy.
  \item[\texttt{norm\_obs=False}] Observation normalisation is disabled because
    all 23 features are already bounded: ownship features in $[-1,1]$ and
    intruder features in $[0,1]$ or $[-1,1]$.
  \item[\texttt{norm\_reward=True}] Reward normalisation is enabled: the running
    mean and variance of the per-environment rewards are tracked by
    \texttt{VecNormalize}, and the reward passed to the value function is
    $r / \sqrt{\hat\sigma^2_r + \epsilon}$.
\end{description}

% -----------------------------------------------------------------------------
\subsection{Callbacks}
% -----------------------------------------------------------------------------

\begin{lstlisting}[caption={Three callbacks attached to the training run.}]
class EpisodeStatsCallback(BaseCallback):
    def _on_step(self):
        for info in self.locals.get('infos', []):
            if 'mean_episode_reward' not in info:
                continue
            self.logger.record('episode/mean_reward', info['mean_episode_reward'])
            self.logger.record('episode/los_steps',   info['ep_los_steps'])
            self.logger.record('episode/length',       info['ep_length'])
            dist  = info.get('action_distribution', [])
            total = max(sum(dist), 1)
            for label, count in zip(self._LABELS, dist):
                self.logger.record(f'actions/{label}', count / total)
        return True
\end{lstlisting}

The three callbacks and their roles:

\begin{description}[noitemsep]
  \item[\texttt{EpisodeStatsCallback}] Reads the \texttt{info} dict populated
    by \texttt{step()} at truncation and logs per-episode metrics to TensorBoard:
    mean reward, LoS step fraction, episode length, and the normalised frequency
    of each action.
    The action histogram is the primary diagnostic for whether the agent has
    learned the turn $\to$ hold $\to$ fly-direct pattern.
  \item[\texttt{CheckpointCallback}] Saves a \texttt{.zip} model checkpoint
    (and the \texttt{VecNormalize} statistics as \texttt{\_vecnorm.pkl}) every
    100\,000 timesteps.
    Both files are needed to resume training or run evaluation.
  \item[\texttt{ProgressCallback}] Prints step count, percentage, elapsed time,
    and throughput (steps/s) to stdout every 10\,000 timesteps.
\end{description}

% -----------------------------------------------------------------------------
\subsection{Training loop and multi-seed mode}
% -----------------------------------------------------------------------------

\begin{lstlisting}[caption={The \texttt{train()} function and \texttt{main()} entry point.}]
def train(seed):
    run_name = f"dalmau_8act_obs23_seed{seed}_{datetime.now().strftime('%Y%m%d_%H%M%S')}"
    run_dir  = os.path.join(RUNS_ROOT, run_name)

    venv = make_vec_env(AirspaceEnv, n_envs=N_ENVS, vec_env_cls=SubprocVecEnv, seed=seed)
    env  = VecNormalize(VecMonitor(venv), norm_obs=False, norm_reward=True, gamma=0.99)
    model = PPO('MlpPolicy', env, seed=seed, tensorboard_log=tb_dir, **PPO_KWARGS)

    model.learn(TOTAL_TIMESTEPS, callback=callbacks, ...)

    model.save(os.path.join(run_dir, 'final_model'))
    env.save(os.path.join(run_dir, 'final_vecnorm.pkl'))

def main():
    if args.multi:
        seeds = [random.randint(1, 99_999) for _ in range(N_RUNS)]
        procs = [subprocess.Popen([sys.executable, __file__, '--seed', str(s)])
                 for s in seeds]
        for p in procs: p.wait()
    else:
        train(args.seed or random.randint(0, 99_999))
\end{lstlisting}

The \texttt{--multi} flag launches \texttt{N\_RUNS = 3} independent training
processes as subprocesses, each with a different random seed.
This is a lightweight form of parallel seed training: the subprocesses run
independently on separate CPU cores and share no state.
Each produces its own checkpoint directory under \texttt{Runs\_saved/v3/}.

The environment wrapper stack is:
\[
  \underbrace{\texttt{AirspaceEnv}}_{\text{raw env}}
  \xrightarrow{\texttt{SubprocVecEnv}}
  \underbrace{\texttt{VecEnv}}_{\text{4 parallel}}
  \xrightarrow{\texttt{VecMonitor}}
  \underbrace{\texttt{VecMonitor}}_{\text{episode logging}}
  \xrightarrow{\texttt{VecNormalize}}
  \underbrace{\texttt{VecNormalize}}_{\text{reward normalisation}}
  \xrightarrow{\texttt{PPO}} \text{policy}
\]

% =============================================================================
\section{End-to-end data flow}
% =============================================================================

The following summarises the full data flow for a single RL step:

\begin{enumerate}
  \item \textbf{Agent selects action} $a \in \{0,\ldots,7\}$ based on current
        observation $\mathbf{o}_t$.
  \item \textbf{\texttt{step(a)} is called.}
  \item Pending aircraft re-enter the sector if their countdown has reached 1.
  \item The action is applied to the focus aircraft via a BlueSky \texttt{HDG}
        command; for turn actions $a \in \{0\text{--}2, 4\text{--}6\}$ the
        commanded heading is updated as
        $\psi_\text{cmd} \leftarrow (\psi_\text{cmd} + \delta_a) \bmod 360^\circ$.
  \item BlueSky advances 10 physics sub-steps ($5\,\text{s}$ simulated time);
        aircraft turn smoothly toward their commanded headings.
  \item Geometric LoS check: any pair within 5\,NM sets
        \texttt{\_los\_this\_step = True}.
  \item Exited aircraft are removed; their slots are queued for respawn.
  \item The urgency matrix is rebuilt from the current aircraft positions and
        the new focus aircraft is selected.
  \item The reward $r_t = r_\text{LoS} + r_\text{conflict} + r_\text{drift}
        + r_\text{work}$ is computed for the acting aircraft.
  \item The observation $\mathbf{o}_{t+1}$ is constructed for the new focus
        aircraft and returned together with $r_t$.
\end{enumerate}

\end{document}
